\chapter{Discussion and Development}
\label{ch:discussion}

This chapter discusses the results of the practical work presented in
Chapter~\ref{ch:development}. It compares the developed bot with similar
existing solutions, proposes a set of improvements for future versions, and
evaluates the feasibility of the system in terms of cost and scalability.
The chapter ends with sub-conclusions that lead to the final Conclusions of
the thesis.

\section{Extended Discussion}
\label{sec:extended-discussion}

The results from Chapter~\ref{ch:development} show that even a small
Telegram bot with a manually collected dataset can solve a real daily
problem for Fergana residents. This section reflects on these results from
three different angles: comparison with related work, the role of the
manual data approach, and the lessons learned during the development.

\subsection{Comparison with Related Work}
\label{subsec:comparison}

Chapter~\ref{ch:literature} reviewed several public transport information
tools, including mobile applications, web services, map platforms (such as
Yandex Maps), and messaging bots. Table~\ref{tab:comparison-discussion}
compares these tools with the developed Fergana Transport Bot from a
practical point of view.

\begin{table}[h]
\centering
\caption{Practical comparison of public transport information tools}
\label{tab:comparison-discussion}
\begin{tabular}{lccc}
\hline
\textbf{Feature} & \textbf{Mobile app} & \textbf{Yandex Maps} &
\textbf{Fergana Bot} \\
\hline
Installation required & yes & yes & no \\
Works inside Telegram & no & no & yes \\
Full coverage of Fergana routes & rare & partial & yes (six routes) \\
Free for the user & varies & yes & yes \\
Development cost & high & --- & low \\
Update by maintainer & app store & company & direct (SQLite) \\
Offline word-of-mouth replacement & possible & possible & yes \\
\hline
\end{tabular}
\end{table}

The comparison shows that the bot is not the most powerful tool in absolute
terms. A mobile application can offer richer features such as maps and
push notifications. Yandex Maps can provide visual navigation. However,
the bot has clear practical advantages for the specific context of
Fergana: it requires no installation, it covers local routes that Yandex
Maps misses, and it can be updated quickly by the maintainer without
publishing a new release.

\subsection{The Manual Data Approach as a Strength}
\label{subsec:manual-data}

At the start of the project, the manual collection of route data was seen
as a weakness, because no official open data was available. After
completing the development, this approach turned out to be a strength for
three reasons.

First, the manual approach forced the author to verify each stop and each
route number with real users from Fergana. This gave the dataset a high
local accuracy that an automatic import from a public API could not
guarantee.

Second, the structured SQLite database makes future updates easy. New
routes can be added without changing the code. This is important in a
city where bus routes can change without official announcement.

Third, the manual approach is reproducible for other regional cities of
Uzbekistan. A student or a small team can collect data in a similar way
for another city and reuse the same bot code.

\subsection{Lessons Learned}
\label{subsec:lessons}

Several lessons emerged during the development. They are summarized in
Table~\ref{tab:lessons}.

\begin{table}[h]
\centering
\caption{Lessons learned from the development of the bot}
\label{tab:lessons}
\begin{tabular}{p{0.32\textwidth}p{0.6\textwidth}}
\hline
\textbf{Lesson} & \textbf{Explanation} \\
\hline
Examples are essential & Test users could not use \texttt{/poisk} until a
clear example was shown in the help message. \\
Buttons feel safer than commands & Several users asked for a button menu
because typed commands felt risky for them. \\
Data quality is harder than code & Writing the bot was faster than
collecting and cleaning the route data. \\
Async pays off only at scale & With 12 test users, async support was not
needed, but it will be needed when the bot is published. \\
Small scope is a strength & A focused MVP was easier to defend than an
ambitious half-finished system. \\
\hline
\end{tabular}
\end{table}

These lessons match a wider theoretical observation: chat bots can
replace simple mobile applications for information delivery in regional
cities, especially when development resources are limited. \cite{adamopoulou2020chatbots, kamberi2021telegram}

\section{Proposals and Improvements}
\label{sec:proposals}

This section proposes a set of improvements for future versions of the
Fergana Transport Bot. The improvements are ordered by priority. The
proposed extended architecture is shown in
Figure~\ref{fig:extended-arch}.

\begin{figure}[h]
\centering
\begin{tikzpicture}[
    node distance=1.1cm and 0.9cm,
    box/.style={rectangle, draw, minimum width=2.2cm, minimum height=0.9cm,
        align=center, font=\small},
    new/.style={rectangle, draw, dashed, minimum width=2.2cm,
        minimum height=0.9cm, align=center, font=\small, fill=green!8},
    db/.style={cylinder, draw, shape border rotate=90, aspect=0.3,
        minimum width=1.8cm, minimum height=1cm, align=center, font=\small},
    arrow/.style={->, >=Stealth, thick}
]
    \node[box] (user) {User};
    \node[box, right=of user] (tg) {Telegram};
    \node[box, right=of tg] (bot) {Python Bot\\(aiogram)};
    \node[db, right=of bot] (db) {SQLite};

    \node[new, below=of tg] (uz) {Uzbek\\language};
    \node[new, below=of bot] (ai) {AI assistant\\(Grok / GPT)};
    \node[new, below=of db] (admin) {Web admin\\panel};

    \draw[arrow] (user) -- (tg);
    \draw[arrow] (tg) -- (bot);
    \draw[arrow] (bot) -- (db);
    \draw[arrow, dashed] (tg) -- (uz);
    \draw[arrow, dashed] (bot) -- (ai);
    \draw[arrow, dashed] (db) -- (admin);
\end{tikzpicture}
\caption{Extended architecture with proposed improvements (dashed = future)}
\label{fig:extended-arch}
\end{figure}

\subsection{Improvement 1: Uzbek Language Support}
\label{subsec:uzbek}

The most common request from test users was to add an Uzbek language
version of the bot. The current version supports only Russian and English
commands. Many Fergana residents speak Uzbek as the first language. Adding
Uzbek support would increase the audience of the bot significantly.

The technical work for this improvement is small. The texts of the bot
replies can be moved to a separate dictionary file. A new
\texttt{/language} command can let the user select the preferred language.
The database does not need any change for this feature.

\subsection{Improvement 2: Button Menu (Reply Keyboard)}
\label{subsec:buttons}

Three test users asked for a button menu instead of typed commands. The
aiogram library supports reply keyboards out of the box. Adding a keyboard
with buttons such as \textit{Routes}, \textit{Search}, \textit{Taxi}, and
\textit{Help} would make the bot easier to use for people who are not
familiar with command-line style interfaces.

\subsection{Improvement 3: AI Assistant Integration}
\label{subsec:ai}

The current bot understands only fixed commands. If the user writes a
free-form question, the bot replies with an "unknown command" message
(see Figure~\ref{fig:botpoisk} in Chapter~\ref{ch:development}). A
practical improvement is to connect the bot with an AI assistant such as
Grok or GPT.

The integration would work as follows. When the bot receives a message
that does not match any known command, it would send the message to the
AI assistant together with a short context about the available routes
and stops. The AI assistant would then generate a free-form answer. The
answer would be returned to the user inside Telegram.

This improvement was intentionally moved to future work. The reason is
that AI assistant integration adds external dependencies, possible API
costs, and a new layer of testing. In the scope of this bachelor thesis,
a focused MVP without AI was a more realistic goal.

\subsection{Improvement 4: Web Admin Panel}
\label{subsec:admin}

The current version of the bot is updated by editing the SQLite database
directly. This is acceptable for a single maintainer but not for a small
team. A simple web admin panel built with a framework such as Flask or
FastAPI would let several maintainers add new routes, edit stop names,
and approve user suggestions from the \texttt{/dobavit} command.

\subsection{Improvement 5: Push Notifications and Schedules}
\label{subsec:push}

Several test users assumed that the bot would know the schedule of buses.
The current version does not include schedules because no public schedule
data exists in Fergana. If a schedule is collected in the future (for
example, with the help of local transport companies), it can be added to
the database and used for two new features:

\begin{itemize}
    \item A \texttt{/schedule <route>} command that shows the next
        departures.
    \item Push notifications to users who subscribed to a specific route.
\end{itemize}

\subsection{Improvement 6: GPS-Based Features}
\label{subsec:gps}

If transport companies in Fergana publish GPS data of their vehicles in
the future, the bot can use this data to show the live position of a bus
and estimate the waiting time at a stop. This improvement is the most
demanding from a technical point of view but also the most valuable from
the user's point of view.

\section{Evaluation of Feasibility and Efficiency}
\label{sec:feasibility}

This section evaluates the feasibility of the bot at the current scale and
at a larger scale. The evaluation is divided into three parts: cost
analysis, performance scalability, and risks of expansion.

\subsection{Cost Analysis}
\label{subsec:cost}

The current cost of running the bot is zero. The free tier of
PythonAnywhere provides enough resources for up to about 1000 active users
per day. The estimated cost at larger scales is shown in
Table~\ref{tab:cost}. The values are based on the published pricing of
common cloud providers in 2026.

\begin{table}[h]
\centering
\caption{Estimated hosting cost at different scales}
\label{tab:cost}
\begin{tabular}{lll}
\hline
\textbf{Scale} & \textbf{Hosting plan} & \textbf{Estimated cost} \\
\hline
Up to 1\,000 users / day & PythonAnywhere free & USD 0 / month \\
Up to 10\,000 users / day & PythonAnywhere Hacker & USD 5 / month \\
Up to 100\,000 users / day & Small VPS (1 vCPU, 2~GB) & USD 10--20 / month \\
City-level deployment & Managed VPS + backups & USD 30--50 / month \\
\hline
\end{tabular}
\end{table}

For comparison, the development cost of a native mobile application with
the same functionality would be much higher. Even a small team would
spend several thousand US dollars and several months to publish an
application on the Google Play and Apple App Store. The Telegram bot
avoids both the development cost and the store publishing process.

\subsection{Performance Scalability}
\label{subsec:scalability}

The bot is currently designed for a small load. The asynchronous design
of aiogram allows it to handle many users without code changes. However,
SQLite has a known limit at high write rates. If the number of users
grows above a few thousand active users per day, the database should be
migrated from SQLite to PostgreSQL. This migration is straightforward
because the data model is simple and the bot uses standard SQL queries.\cite{kreibich2010sqlite, sqlite_docs}

The Telegram Bot API also has a rate limit of approximately 30~messages
per second to different users. For a single bot, this limit is high
enough for tens of thousands of users per day under normal load.\cite{telegram_botapi}

\subsection{Risks of Expansion}
\label{subsec:risks}

The main risks of expanding the bot are summarized in
Table~\ref{tab:risks}.

\begin{table}[h]
\centering
\caption{Risks of expanding the bot to a larger user base}
\label{tab:risks}
\begin{tabular}{p{0.32\textwidth}p{0.6\textwidth}}
\hline
\textbf{Risk} & \textbf{Mitigation} \\
\hline
Outdated route data & Add a web admin panel and accept verified user
suggestions through \texttt{/dobavit}. \\
Server downtime & Move to paid hosting with backups and a status
monitor. \\
SQLite write bottleneck & Migrate to PostgreSQL when daily active
users exceed a few thousand. \\
AI integration costs & Use rate limiting and a per-user daily quota for
AI requests. \\
Privacy of user data & Store only the Telegram identifier and the
preferred language; nothing else. \\
\hline
\end{tabular}
\end{table}

Each risk has a clear and low-cost mitigation. None of the risks blocks
the expansion of the bot to a larger user base.

\section{Sub-Conclusions}
\label{sec:ch4-subconclusions}

This chapter discussed the results of the practical work, proposed six
improvements for future versions of the bot, and evaluated the
feasibility of the system at different scales.

The comparison with related work showed that the Fergana Transport Bot is
not the most powerful tool, but it is the most practical solution for the
specific context of Fergana city. The manual data collection approach,
seen as a weakness at the start, turned out to be a strength because of
its local accuracy and reproducibility.

The most promising improvement is the integration with an AI assistant
such as Grok or GPT, because it removes the limitation of fixed commands.
The most accessible improvement is the addition of Uzbek language
support, because it requires only small code changes and would expand the
user base significantly.\cite{dale2016return}

The cost analysis showed that the bot can grow from zero cost to a small
monthly fee even at city-level scale. The performance scalability is
acceptable for tens of thousands of users with minimal changes. All
identified risks have low-cost mitigations.

These insights show that the developed system is not only a working MVP
but also a realistic starting point for a larger transport information
service. The achievement of the aim and the validation of the hypothesis
are summarized in the next chapter.